\chapter{검증 결과(7): Stage A 경부하 운전과 모듈 관리·고장}

\section{위상 셰딩의 필요성과 설계}

공진 스위치드 커패시터는 공진전류의 영교차에서 꺼 스위칭 손실을 줄인다\cite{b22}. 영교차 추종 게이트는 전류가 검출 문턱(0.5 A) 아래로 돌아온 뒤 검출 지연(20 ns)만큼 늦게 끄므로, 공진전류 첨두가 작으면 이 문턱·지연이 반주기에서 차지하는 몫이 커져 턴오프 전류/피크가 오른다. 같은 출력을 위상 하나가 맡으면 첨두가 두 배가 되어 비율이 절반 아래로 내려가고, 위상당 고정 손실(출력용량·게이트)도 준다. 다상 변환기에서 경부하 효율을 위해 위상 수를 줄이는 것은 널리 쓰는 방법이다\cite{b23}. 문턱은 운전점 계산(표 \ref{t:e24d})으로 정했다. 11.1–11.4절의 시험은 150 kHz Stage A로 했고, 최종 설계(A25R)의 문턱과 판정은 11.5절에 있다. 표 \ref{t:e24d}의 모듈 손실은 출력용량 손실을 정정하기 전 모델 값이다(13장).

{\footnotesize
\begin{longtable}{L{\dimexpr 0.3019\linewidth-2\tabcolsep\relax}L{\dimexpr 0.1132\linewidth-2\tabcolsep\relax}L{\dimexpr 0.1132\linewidth-2\tabcolsep\relax}L{\dimexpr 0.1132\linewidth-2\tabcolsep\relax}L{\dimexpr 0.1132\linewidth-2\tabcolsep\relax}L{\dimexpr 0.1132\linewidth-2\tabcolsep\relax}L{\dimexpr 0.1321\linewidth-2\tabcolsep\relax}}
\caption{모듈 출력전류별 영전류 스위칭 품질과 손실(E24, 1008 V)}\label{t:e24d}\\
\toprule
\tabhead \textbf{모듈 출력전류} & \textbf{5 A} & \textbf{7.5 A} & \textbf{10 A} & \textbf{15 A} & \textbf{22 A} & \textbf{46.9 A} \\
\midrule
\endfirsthead
\toprule
\tabhead \textbf{모듈 출력전류} & \textbf{5 A} & \textbf{7.5 A} & \textbf{10 A} & \textbf{15 A} & \textbf{22 A} & \textbf{46.9 A} \\
\midrule
\endhead
\bottomrule
\endlastfoot
2위상 턴오프 전류/피크 & 12.3 \% & 7.19 \% & \textbf{5.51 \%} & 3.78 \% & 2.14 \% & 0.46 \% \\
위상 하나 턴오프 전류/피크 & 4.97 \% & 3.27 \% & 1.58 \% & 1.58 \% & 0.11 \% & 0.11 \% \\
위상 하나 첨두 & 8.0 A & 12.1 A & 16.2 A & 24.3 A & 35.9 A & \textbf{76.4 A} \\
모듈 손실 2위상 / 위상 하나 & 31.6 / 15.8 W & 31.4 / 16.4 W & 31.7 / 17.2 W & 32.8 / 19.7 W & 35.2 / 25.1 W & 52.6 / 60.6 W \\
\end{longtable}
}
\srcnote{주: 46.9 A는 충전 정격점(788 V)의 위상 전류로, 이때 위상 하나의 첨두 76.4 A를 위상 정격 첨두로 둔다. 540 V도 거의 같다.}

\begin{itemize}
\item \textbf{멈춤 문턱 15 A(100 µs 지속).} 2위상은 12 A 이상에서 기준 5 \% 안이므로 3 A 여유를 둔다.
\item \textbf{복귀 문턱 22 A 또는 남은 위상의 순시 전류 36 A.} 15–22 A는 히스테리시스 폭이고 두 운전 모두 기준 안이다. 필터한 평균전류는 부하 계단을 늦게 보므로 순시 전류도 본다. 위상 하나의 정상 첨두는 20 A 이하에서 32.6 A 이하라 36 A는 정상 운전에서 잘못 복귀시키지 않는다.
\item \textbf{멈출 때는 도통 중인 구간을 자르지 않는다.} 지금 구간이 자연 턴오프로 끝난 뒤 막으므로 멈추는 순간도 영전류다. 복귀는 바로 한다.
\end{itemize}

확인 과정에서 Stage A 스위칭 모델의 데드타임 처리 오차를 찾아 고쳤다. 데드타임 동안 다이오드가 환류하는 전류가 적분 간격 안에 0으로 가지 않고 남았고, 멈춘 위상에서는 이 잔류 전류가 진동해 재시작 위상의 영교차 검출을 흔들었다. 고친 처리로 회로(TB16·TB17) V\textsubscript{A}와의 최대 차가 0.094 V에서 0.014 V로 줄었다. 이전 처리는 턴오프 비율을 높게(보수 쪽으로) 계산했으므로 E1–E23의 판정은 바뀌지 않고, E24 이후의 기준값은 고친 처리로 계산했다.

\section{위상 셰딩의 결과}

셰딩은 2위상 정상 운전에 없던 과도(멈추는 순간, 위상 하나의 부하 계단, 멈춰 있던 위상의 재시작, 문턱 근처의 반복 전환)를 만든다. 그래서 정상 운전점과 함께 전환 과도를 시험 전에 정한 기준(모든 턴오프 |i\textsubscript{off}| ≤ max(5 \% × 반주기 첨두, 1.0 A), 위상 첨두 ≤ 76.4 A, 반복 전환 없음)으로 판정했다(표 \ref{t:e24}).

{\footnotesize
\begin{longtable}{L{\dimexpr 0.2885\linewidth-2\tabcolsep\relax}L{\dimexpr 0.5962\linewidth-2\tabcolsep\relax}L{\dimexpr 0.1154\linewidth-2\tabcolsep\relax}}
\caption{Stage A 경부하 위상 셰딩의 판정(E24, TB18, TB19)}\label{t:e24}\\
\toprule
\tabhead \textbf{시험} & \textbf{결과} & \textbf{판정} \\
\midrule
\endfirsthead
\toprule
\tabhead \textbf{시험} & \textbf{결과} & \textbf{판정} \\
\midrule
\endhead
\bottomrule
\endlastfoot
E24 정상 운전점(1008·540 V × −10–75 A) & 7.5 A 이상 턴오프 전류/피크 ≤ 3.27 \%(셰딩 없는 10 A 5.51 \%). 7.5 A 미만은 턴오프 전류 ≤ 0.48 A & 26/26 \\
E24 과도(계단·램프·역방향·충전기 주입 10사례) & 턴오프 위반 0, 반복 전환 없음. V\textsubscript{A} 노드 직접 계단에서 위상 첨두 72.9 A(셰딩 없는 60.5 A). 사례 E(모듈당 충전 램프 요구의 4배)는 셰딩 없이도 82.5 A, 셰딩 93.6 A로 정격 초과 & 9/10 \\
TB18 회로(R2 버스 뒤 부하 계단·램프, 충전기 주입, 역방향 5사례) & 턴오프 3,547번 위반 0(최대 0.77 A), 위상 첨두 ≤ 52.5 A(셰딩 없는 회로보다 ≤ 5.1 A 큼). E24와 전환 시각 2.3 µs, 첨두 0.23 A, V\textsubscript{A} 0.005 V 안 & 5/5 \\
TB19 회로(TB17 + 셰딩 게이트) & 1008 V 강등 뒤 10 A의 턴오프 전류/피크 4.44 → 1.31 \%. 턴오프 17,561번 위반 0, 위상 첨두 ≤ 48.6 A. 멈춤·복귀(3·2번, 4·4번)가 모두 실제 전류 변화를 따랐고 취소 0. 강등·48 V·Stage C 결과는 TB17과 같음 & 2/2 \\
\end{longtable}
}
\srcnote{주: 7.5 A 미만에서는 턴오프 전류/피크 비율이 뜻이 없어 절대값(≤ 1.0 A, 검출 문턱의 두 배)으로 판정했다.}

\begin{figure}[H]
\centering
\includegraphics[width=15.5cm]{\repo/validation/2026-10-26-stage-a-shedding/E24/E24_stage_a_shedding.png}
\caption{Stage A 경부하 위상 셰딩(E24)}\label{f:e24}
\srcnote{주: (a) 모듈 출력전류별 턴오프 전류/피크, (b) 모듈 손실, (c) 54 → 10 → 54 A 계단의 위상 전류와 셰딩 모드, (d) 사례 E(모듈당 충전 램프 요구의 4배)의 위상 전류.}
\end{figure}

불통과 1건(사례 E)은 셰딩이 없어도 위상 정격을 넘는 조건이다. 셰딩은 복귀 지연만큼 첨두를 11.1 A 더했다. 150 kHz 설계의 충전전류 증가 램프 요구(전체 3.75 A/µs, 모듈당 0.94 A/µs, 8.3절)에서는 위상 첨두가 52.6 A로 셰딩 없는 경우(53.1 A)와 같았다. 따라서 충전기 쪽 램프 제한이 위상 정격도 지킨다. TB19에서는 강등 중 클램프 유지 구간(약 3 A)의 반주기별 턴오프 비율도 약 24 \%에서 약 10 \%로 내려갔다(턴오프 전류 ≤ 0.43 A). 이 시험은 모듈 하나만 계산했다. 모듈 4개(위상 8개)에서 모듈 셰딩과 위상 셰딩의 조합, 부품 공차, 멈춘 모듈의 재시작은 다음 절에서 다뤘다.

\section{4모듈 관리기와 부품 공차(E28, TB22)}

실제 Stage A는 모듈 4개이고, 운전점마다 필요한 모듈 수가 다르다. 주행 보조부하는 수–100 A(모듈 1–2개), 400–500 V급 충전은 375 A에서 보조부하를 뺀 값(모듈 4개), R2-W460 강등 뒤는 약 10 A(모듈 1개의 위상 하나)다. 모듈 넷을 늘 돌리면 경부하에서 영전류 스위칭 비율이 나빠지고 위상당 고정 손실이 여덟 번 든다. 그래서 동작 모듈 수를 정하는 관리기를 설계했다.

\begin{itemize}
\item \textbf{추가(바로).} 전체 출력전류 > 동작 수 × 70 A, 위상 순시 전류 > 66 A, 충전 요청 가운데 하나면 모듈을 더한다.
\item \textbf{줄임(지속 뒤).} $\lceil$전체 전류/60 A$\rceil$ < 동작 수가 1 ms 이어지면 줄인다. 줄인 모듈은 지금 구간을 자연 턴오프로 끝내고 멈춘다. 줄임 문턱 60·120·180 A는 모듈 N개와 N+1개의 손실이 같아지는 전류보다 낮게 두어, 줄인 쪽의 손실이 작거나 같게 했다.
\item \textbf{충전 시작.} 충전 요청이 오면 충전기 전류가 오르기 전에 모듈을 모두 켠다(사전 기동).
\item \textbf{동작 모듈.} 모듈마다 11.1절의 위상 셰딩을 한다.
\end{itemize}

스위칭 모델(4모듈 8위상, R2 버스, V\textsubscript{A} 노드 충전기)로 공칭과 공차 표본 3개(C 5 \%, L 10 \%, R 10 \%)를 계산했다. 판정 기준은 시험 전에 정했고, 1·2차 실행의 불통과로 설계를 네 군데 고친 뒤 3차 실행으로 판정했다.

\begin{itemize}
\item \textbf{정상 운전점: 공칭 33/36, 공차 99/108.} 주행 24점·충전 12점에서 동작 모듈 수와 위상 수가 규칙과 같았고, 턴오프 비율 ≤ 4.1 \%, 턴오프 전류 ≤ 0.91 A였다. 불통과는 정격 충전점(375 A)의 첨두 판정뿐이다. 판정 창(0.1 ms 이후 전체)에 공진전류 0에서 시작하는 시작 과도(82.1 A)가 들어가기 때문이다. 정상 창의 첨두는 모두 정격 이하였고, 판정 창은 바꾸지 않았다.
\item \textbf{과도: 공칭 6/7, 공차 21/21.} 부하 계단·램프, 충전 시작·종료(사전 기동), 1008 V 경부하에서 턴오프 위반이 없었고 멈추기 취소도 없었다. 불통과는 사전 기동이 없는 충전(위상 첨두 84.2 A)으로, 이 결과로 사전 기동을 요구사항에 넣었다.
\item \textbf{회로(TB22): 사전 기동이 있는 세 사례 통과, 사전 기동 없는 사례 불통과(모델과 같음).} 모듈 수 전환 시각은 모델과 2 µs 안, 위상 첨두는 3.3 A 안이었다.
\item \textbf{부품 공차.} 공차 표본의 정격 충전점 위상 첨두는 공칭보다 10–12 \% 컸고, 위상 평균전류 불균형은 최대 8.5 \%였다. 이 값은 부품 정격 표(16장)의 입력으로 넘겼다.
\item \textbf{멈춘 모듈의 플라잉 커패시터 전압.} 재시작은 C\textsubscript{fly} 전압 오차가 V\textsubscript{bat}/2 대비 −10–0 V일 때만 판정을 지켰다(+2 V에서 남은 모듈 첨두 77.1 A, −15 V에서 재시작 첨두 106.9 A). 정지 모듈의 균형 회로가 C\textsubscript{fly}를 V\textsubscript{bat}/2보다 0–10 V 낮게 유지해야 한다. 이 요구는 150 kHz 설계에서만 확인했다.
\item \textbf{게이트 변경.} 무장하지 못한 반주기(첨두 < 2.5 A)는 부호 없이 |i| < 0.5 A가 검출 지연만큼 이어질 때 끄도록 바꿨다(12.3절의 방향 반전 처리를 대체). 바꾼 게이트로 E27 사례 55개와 E24 판정을 다시 확인해 결과가 같았다.
\end{itemize}

\section{모듈 고장 처리(E30, TB24)}

E28의 관리기는 고장 난 모듈을 후보에서 빼는 처리가 없었다. 주행 중 유일한 동작 모듈이 고장 나면 관리기가 그 모듈을 계속 ``동작''으로 세어 대신할 모듈이 시작하지 않는다. 충전 중 고장이면 충전기가 전류를 줄일 때까지 수십 ms 이상 남은 3모듈이 375 A를 받는다. 그래서 고장 처리를 넣었다. 고장 간격에 고장 모듈 위상의 네 스위치를 끄고(남은 공진전류는 바디 다이오드로 소멸), 건전 모듈 목록에서 영구히 빼며, 동작 모듈이 0이 되면 다음 건전 모듈을 바로 두 위상으로 시작한다.

주행·충전 중 고장 8사례를 공칭·공차 표본으로 판정했다. 기준은 턴오프 |i\textsubscript{off}| ≤ max(5 \% × 반주기 첨두, 1.0 A), 위상 첨두 주행 ≤ P\textsubscript{rated}·충전 ≤ 1.5 × P\textsubscript{rated}(표본별), V\textsubscript{A} 고장 직전 ±10 V, 끝 동작 수가 기대값인 것이다. \textbf{모델 32/32, 회로(TB24) 4/4로 통과했다.} 주행 중 유일한 동작 모듈이 고장 나면 고장 간격에 다음 모듈이 시작했고(대체 지연 0 µs) V\textsubscript{A} 변화는 0.2 V 안이었다. 충전 중 고장에서는 남은 3모듈의 위상 첨두가 103.5–114.5 A(한도 114.7–128.3 A)였고, 이 값이 Stage A 소자·인덕터·플라잉 커패시터의 정격 조건이 되었다(16장). 회로는 위상 첨두가 모델과 0.9 A 안, V\textsubscript{A}가 0.05 V 안이었다.

\section{최종 설계(A25R)의 문턱과 판정}

Stage A를 25 kHz 설계 A25R로 바꾸면서(13장) 문턱과 시간 상수를 다시 정하고, 같은 판정 기준으로 다시 확인했다(E32, E35, TB27). 결과는 표 \ref{t:a25r_mgr}와 같다.

{\footnotesize
\begin{longtable}{L{\dimexpr 0.26\linewidth-2\tabcolsep\relax}L{\dimexpr 0.30\linewidth-2\tabcolsep\relax}L{\dimexpr 0.44\linewidth-2\tabcolsep\relax}}
\caption{Stage A 경부하·모듈 관리의 150 kHz 설계 값과 최종 설계 값}\label{t:a25r_mgr}\\
\toprule
\tabhead \textbf{항목} & \textbf{150 kHz 설계} & \textbf{A25R(최종)} \\
\midrule
\endfirsthead
\toprule
\tabhead \textbf{항목} & \textbf{150 kHz 설계} & \textbf{A25R(최종)} \\
\midrule
\endhead
\bottomrule
\endlastfoot
위상 셰딩(멈춤 / 복귀 / 순시) & 15 / 22 / 36 A & 20 / 30 / 53 A(두 위상 7.5–12.5 A의 턴오프 비율 때문) \\
모듈 관리기(추가 / 줄임 / 순시) & 70 / 60 A/모듈, 66 A & 같음(손실 교차 89·156·221 A ≥ 줄임 60·120·180 A) \\
시간 상수(필터, 멈춤 지속, 모듈 필터, 잠금, 유지) & 2 µs, 100 µs, 20 µs, 20 µs, 100 µs & 6배: 12 µs, 600 µs, 120 µs, 120 µs, 600 µs(모듈 줄임 지속 1 ms는 그대로) \\
정격 충전점 위상 첨두 P\textsubscript{rated} & 76.4–76.5 A & 74.5 A(공차 표본 86.5–102.8 A) \\
충전 전 사전 기동·충전전류 램프 & ≥ 100 µs & ≥ 600 µs \\
정상 운전점(공칭 / 공차 표본) & 33/36 / 99/108 & 33/36 / 36/36(공칭 불통과는 정격 충전점 시작 과도 76.7 A, 마지막 10주기 74.3 A) \\
요구 안 과도(공칭 / 표본) & 6/7 / 21/21 & 6/6 / 6/6 \\
모듈 고장(검출 즉시 해당 모듈 게이트 차단) & 32/32 & 16/16(3모듈 충전 첨두 최대 105.5 A; 후보 C\textsubscript{fly} 45.0·42.75 µF에서 106.86 A) \\
회로 & TB22 3/4, TB24 4/4 & TB27 4/4(첨두 차 ≤ 2.2 A, V\textsubscript{A} 차 ≤ 0.06 V), TB28 4/4 \\
\end{longtable}
}
\srcnote{주: A25R의 셰딩 문턱 20 / 30 / 53 A는 지정한 모듈·부하 이력에서 정한 제어값이다. 충전 시작은 보조부하 15·25 kW에서 확인했고, 640 V·보조부하 1 kW에서는 같은 기동·램프로 78.74 A가 되어 공칭 한도 74.53 A를 넘었다(E46). 150 kHz 설계에서 확인한 프리차지·정지 모듈 재기동은 최종 C\textsubscript{A} 2.4 mF의 냉기동 근거로 쓰지 않는다(최종 설계로 다시 확인하지 않음). 모듈 고장은 고장 표지를 받는 즉시 해당 모듈의 모든 게이트를 끄는 차단 가능한 고장이며, 강제 차단 순간은 정상 턴오프 평가와 분리했다(단락 소자·게이트 고착·검출 지연은 미확인). 부하 계단은 요구사항의 계약(Tier-1 대표 시험 ±10 kW)으로 판정했다. 150 kHz부터 써 온 20 kW 순간 계단은 계약 밖으로 정보로만 계산했고 A25R에서 불통과였다(첨두 87–88 A). 자료: 검증 근거 E24, E28, E30, E32, E35.}

시간 상수를 6배로 늘린 것은 공진 주기가 6배 길어졌기 때문이다. 셰딩 문턱을 올린 것은 A25R의 두 위상 운전이 모듈 출력 7.5–12.5 A에서 턴오프 비율 8.7–5.0 \%로 기준(5 \%)을 넘었기 때문이다. 영교차 문턱(0.5 A)으로 끄는데 반주기 첨두가 작아서다. 위상 하나 운전은 7.5 A 이상에서 4.1 \% 이하다(E32). 멈춘 모듈의 플라잉 커패시터 전압 요구는 A25R에서 다시 확인하지 않았다(18장).
